\documentclass[12pt,a4paper]{article}
\usepackage[utf8]{inputenc}
\usepackage[turkish]{babel}
\usepackage{amsmath, amssymb}
\usepackage{geometry}
\geometry{margin=2.5cm}

\title{Görüntü İşleme - Matematiksel Formüller}
\author{Görüntü İşleme Dersi}
\date{}

\begin{document}
\maketitle

\section{Görüntü Yükleme ve Temel Kavramlar}

\textbf{Dijital Görüntü Matrisi:}
\begin{equation}
I(x,y) \in \mathbb{R}^{M \times N \times C}
\end{equation}

Burada:
\begin{itemize}
    \item $M$ : Yükseklik (satır sayısı)
    \item $N$ : Genişlik (sütun sayısı)
    \item $C$ : Kanal sayısı (Gri=1, RGB=3)
\end{itemize}

\section{Örnekleme (Sampling)}

\textbf{Nyquist-Shannon Örnekleme Teoremi:}
\begin{equation}
f_s \geq 2 \cdot f_{max}
\end{equation}

\textbf{Yeniden Örnekleme (Downsampling):}
\begin{equation}
I_{down}(x,y) = I(x \cdot k, y \cdot k)
\end{equation}

\textbf{Piksel Sayısı:}
\begin{equation}
N_{pixels} = M \times N = \frac{M_{orig}}{k} \times \frac{N_{orig}}{k}
\end{equation}

\section{Kuantalama (Quantization)}

\textbf{Bit Derinliği ve Seviye Sayısı:}
\begin{equation}
L = 2^b
\end{equation}

\textbf{Kuantalama Formülü:}
\begin{equation}
Q(I) = \left\lfloor \frac{I}{step} \right\rfloor \times step
\end{equation}

\begin{equation}
step = \frac{256}{L} = \frac{256}{2^b}
\end{equation}

\textbf{Kuantalama Hatası:}
\begin{equation}
e_q = I - Q(I)
\end{equation}

\textbf{Ortalama Kare Hata (MSE):}
\begin{equation}
MSE = \frac{1}{MN} \sum_{x=0}^{M-1} \sum_{y=0}^{N-1} [I(x,y) - Q(I(x,y))]^2
\end{equation}

\section{Renk Uzayları}

\subsection{RGB (Red, Green, Blue)}
\begin{equation}
\vec{p} = \begin{bmatrix} R \\ G \\ B \end{bmatrix}, \quad R,G,B \in [0, 255]
\end{equation}

\subsection{Grayscale (Gri Tonlama)}

\textbf{Luminosity Yöntemi (ITU-R BT.601):}
\begin{equation}
Y = 0.299 \cdot R + 0.587 \cdot G + 0.114 \cdot B
\end{equation}

\textbf{Ortalama Yöntemi:}
\begin{equation}
Y = \frac{R + G + B}{3}
\end{equation}

\subsection{HSV (Hue, Saturation, Value)}

\textbf{RGB $\rightarrow$ HSV Dönüşümü:}
\begin{equation}
V = \max(R, G, B)
\end{equation}

\begin{equation}
S = \begin{cases} 
0 & \text{if } V = 0 \\ 
\frac{V - \min(R,G,B)}{V} & \text{otherwise} 
\end{cases}
\end{equation}

\begin{equation}
H = \begin{cases} 
0 & \text{if } S = 0 \\
60^{\circ} \times \frac{G-B}{V-min} & \text{if } V = R \\
60^{\circ} \times \left(2 + \frac{B-R}{V-min}\right) & \text{if } V = G \\
60^{\circ} \times \left(4 + \frac{R-G}{V-min}\right) & \text{if } V = B
\end{cases}
\end{equation}

\subsection{YCbCr (Luminance, Chrominance)}
\begin{equation}
\begin{bmatrix} Y \\ Cb \\ Cr \end{bmatrix} = 
\begin{bmatrix} 0.299 & 0.587 & 0.114 \\ -0.169 & -0.331 & 0.500 \\ 0.500 & -0.419 & -0.081 \end{bmatrix}
\begin{bmatrix} R \\ G \\ B \end{bmatrix} +
\begin{bmatrix} 0 \\ 128 \\ 128 \end{bmatrix}
\end{equation}

\section{Enhancement - Histogram Eşitleme}

\begin{equation}
s_k = (L-1) \sum_{j=0}^{k} p_r(r_j) = (L-1) \sum_{j=0}^{k} \frac{n_j}{MN}
\end{equation}

\section{Restorasyon - Gaussian Filtre}

\textbf{2D Gaussian Fonksiyonu:}
\begin{equation}
G(x,y) = \frac{1}{2\pi\sigma^2} e^{-\frac{x^2 + y^2}{2\sigma^2}}
\end{equation}

\section{Morfolojik İşlemler}

\textbf{Dilation (Genişletme):}
\begin{equation}
A \oplus B = \{z : (\hat{B})_z \cap A \neq \emptyset\}
\end{equation}

\textbf{Erosion (Aşındırma):}
\begin{equation}
A \ominus B = \{z : (B)_z \subseteq A\}
\end{equation}

\textbf{Opening:}
\begin{equation}
A \circ B = (A \ominus B) \oplus B
\end{equation}

\textbf{Closing:}
\begin{equation}
A \bullet B = (A \oplus B) \ominus B
\end{equation}

\section{Thresholding (Eşikleme)}

\textbf{Global Threshold:}
\begin{equation}
g(x,y) = \begin{cases} 255 & \text{if } f(x,y) > T \\ 0 & \text{otherwise} \end{cases}
\end{equation}

\textbf{Otsu Threshold:}
\begin{equation}
\sigma_B^2(t) = \omega_0(t)\omega_1(t)[\mu_0(t) - \mu_1(t)]^2
\end{equation}

\begin{equation}
t^* = \arg\max_{t} \sigma_B^2(t)
\end{equation}

\section{Edge Detection (Kenar Tespiti)}

\textbf{Sobel X:}
\begin{equation}
G_x = \begin{bmatrix} -1 & 0 & +1 \\ -2 & 0 & +2 \\ -1 & 0 & +1 \end{bmatrix} * I
\end{equation}

\textbf{Sobel Y:}
\begin{equation}
G_y = \begin{bmatrix} -1 & -2 & -1 \\ 0 & 0 & 0 \\ +1 & +2 & +1 \end{bmatrix} * I
\end{equation}

\textbf{Gradyan Büyüklüğü:}
\begin{equation}
G = \sqrt{G_x^2 + G_y^2}
\end{equation}

\textbf{Gradyan Yönü:}
\begin{equation}
\theta = \arctan\left(\frac{G_y}{G_x}\right)
\end{equation}

\section{K-Means Segmentasyon}

\textbf{Amaç Fonksiyonu:}
\begin{equation}
J = \sum_{i=1}^{K} \sum_{x \in C_i} \|x - \mu_i\|^2
\end{equation}

\textbf{Merkez Güncelleme:}
\begin{equation}
\mu_i = \frac{1}{|C_i|} \sum_{x \in C_i} x
\end{equation}

\section{Connected Components}

\textbf{8-Bağlantılı Komşuluk:}
\begin{equation}
N_8(p) = \{q : \|p - q\|_\infty \leq 1\}
\end{equation}

\section{Emboss (Kabartma) Efekti}

\begin{equation}
K_{emboss} = \begin{bmatrix} -2 & -1 & 0 \\ -1 & 1 & 1 \\ 0 & 1 & 2 \end{bmatrix}
\end{equation}

\section{Fourier Dönüşümü}

\textbf{2D Ayrık Fourier Dönüşümü:}
\begin{equation}
F(u,v) = \sum_{x=0}^{M-1} \sum_{y=0}^{N-1} f(x,y) \cdot e^{-j2\pi\left(\frac{ux}{M} + \frac{vy}{N}\right)}
\end{equation}

\textbf{Ters Fourier Dönüşümü:}
\begin{equation}
f(x,y) = \frac{1}{MN} \sum_{u=0}^{M-1} \sum_{v=0}^{N-1} F(u,v) \cdot e^{j2\pi\left(\frac{ux}{M} + \frac{vy}{N}\right)}
\end{equation}

\section{Hough Dönüşümü}

\textbf{Çizgi Tespiti (Polar Form):}
\begin{equation}
\rho = x \cos\theta + y \sin\theta
\end{equation}

\textbf{Daire Tespiti:}
\begin{equation}
(x - a)^2 + (y - b)^2 = r^2
\end{equation}

\section{Konvolüsyon}

\textbf{2D Konvolüsyon:}
\begin{equation}
(I * K)(x,y) = \sum_{i=-k}^{k} \sum_{j=-k}^{k} I(x-i, y-j) \cdot K(i,j)
\end{equation}

\end{document}
